\documentclass[12pt]{article}
\usepackage[utf8]{inputenc}
\usepackage[T2A]{fontenc}
\usepackage[russian]{babel}
\usepackage{geometry}
\usepackage{graphicx}
\usepackage{float}
\usepackage{booktabs}
\usepackage{array}
\usepackage{amsmath}
\usepackage{caption}
\usepackage{tabularx}
\usepackage{ragged2e}
\usepackage{makecell}
\geometry{margin=2.5cm}
\emergencystretch=3em
\captionsetup{font=small, labelfont=bf}
\newcolumntype{Y}{>{\RaggedRight\arraybackslash}X}

\title{Анализ и прогноз числа зарегистрированных браков в Москве и Архангельской области}
\author{Ковалёва Мария}
\date{30 апреля 2026 года}

\begin{document}
\maketitle

\section{Описательный анализ}
В работе анализируются месячные временные ряды числа зарегистрированных браков.
Рассматриваются два региона: Москва и Архангельская область. Такой выбор
позволяет сопоставить крупнейший столичный мегаполис и северный регион,
отличающиеся по климату, степени урбанизации и демографической структуре.
Именно этот контраст делает сравнение содержательно полезным при изучении уровня,
вариативности и сезонности брачности.

\subsection{Сравнение уровней двух регионов}
На рисунке~\ref{fig:comparison_ts} представлена совместная динамика числа браков
в Москве и Архангельской области.

\begin{figure}[H]
\centering
\includegraphics[width=0.9\textwidth]{figures/regions_comparison_timeseries.png}
\caption{Сравнение динамики числа зарегистрированных браков по двум регионам}
\label{fig:comparison_ts}
\end{figure}

Москва на всем интервале наблюдений заметно превосходит Архангельскую область по
абсолютному уровню показателя. Это естественно связано с большей численностью
населения, высокой концентрацией городских взаимодействий и более масштабным
рынком браков. Архангельская область, напротив, демонстрирует существенно меньший
масштаб и более выраженное относительное снижение к концу выборки, что можно
связать с демографическим сжатием и менее интенсивным миграционным притоком.

\subsection{Динамика Москвы}
На рисунке~\ref{fig:moscow_ts} показан временной ряд числа браков в Москве.

\begin{figure}[H]
\centering
\includegraphics[width=0.85\textwidth]{figures/moscow_timeseries.png}
\caption{Число зарегистрированных браков: Москва}
\label{fig:moscow_ts}
\end{figure}

Московский ряд колеблется вокруг высокого уровня и сохраняет устойчивую
внутригодовую повторяемость. Это соответствует специфике крупной агломерации,
где высокая плотность населения и более молодая возрастная структура поддерживают
значительный объем регистраций на протяжении всего года. Долгосрочный тренд в
выборке выражен умеренно: основная динамика определяется сезонной, а не
направленной компонентой.

\subsection{Динамика Архангельской области}
На рисунке~\ref{fig:ark_ts} представлен ряд числа браков в Архангельской области.

\begin{figure}[H]
\centering
\includegraphics[width=0.85\textwidth]{figures/arkhangelsk_oblast_timeseries.png}
\caption{Число зарегистрированных браков: Архангельская область}
\label{fig:ark_ts}
\end{figure}

Ряд Архангельской области имеет тот же общий сезонный рисунок, но намного меньший
уровень и более заметную относительную амплитуду колебаний. Для северного региона
это выглядит правдоподобно: меньшая численность населения делает показатель
чувствительнее к демографическим и календарным сдвигам. Кроме того, более
суровые климатические условия усиливают концентрацию свадебных регистраций в
теплые месяцы.

\subsection{Описательные статистики}
В таблице~\ref{tab:desc} приведены основные описательные статистики для двух
регионов.

\begin{table}[H]
\centering
\caption{Описательные статистики}
\label{tab:desc}
\small
\begin{tabular}{lrrrr}
\toprule
Регион & Среднее & Ст. откл. & Мин. & Макс. \\
\midrule
Архангельская область & 834.26 & 319.30 & 280 & 1620 \\
Москва & 7619.53 & 2819.36 & 3100 & 13915 \\
\bottomrule
\end{tabular}
\end{table}

Среднее число браков в Москве примерно в девять раз выше, чем в Архангельской
области. Это отражает не только разницу в размере населения, но и различия в
урбанизации, интенсивности миграции и возрастной структуре. Более высокий
абсолютный разброс в Москве также естественен для крупного мегаполиса с большим
числом регистраций.

\section{Сезонность}
\subsection{Средние значения по месяцам}
На рисунке~\ref{fig:seasonal_avg} показаны средние месячные значения числа
браков для двух регионов.

\begin{figure}[H]
\centering
\includegraphics[width=0.85\textwidth]{figures/seasonal_average_by_month.png}
\caption{Среднее число браков по месяцам}
\label{fig:seasonal_avg}
\end{figure}

В обоих регионах наблюдается выраженная годовая сезонная волна. Максимум в
среднем приходится на лето и начало осени, особенно на август, тогда как зимние
месяцы характеризуются более низким числом регистраций. Такой профиль согласуется
с календарной логикой брачного поведения: теплый сезон удобнее для проведения
церемоний и сопутствующих мероприятий.

\subsection{Распределение по месяцам: Москва}
На рисунке~\ref{fig:moscow_box} представлено распределение числа браков по
месяцам для Москвы.

\begin{figure}[H]
\centering
\includegraphics[width=0.75\textwidth]{figures/moscow_monthly_boxplot.png}
\caption{Распределение числа браков по месяцам: Москва}
\label{fig:moscow_box}
\end{figure}

Для Москвы сезонность не является случайным эпизодом: летние месяцы устойчиво
имеют более высокие медианы и квартильные значения. Это естественно для
столичного региона, где большой городской рынок браков поддерживает высокий
уровень регистраций даже вне пикового сезона, но в летне-осенние месяцы спрос
дополнительно усиливается. Следовательно, сезонная компонента в Москве выражена
сильно, хотя относительная вариация ниже, чем в северном регионе.

\subsection{Распределение по месяцам: Архангельская область}
На рисунке~\ref{fig:ark_box} показано распределение числа браков по месяцам для
Архангельской области.

\begin{figure}[H]
\centering
\includegraphics[width=0.75\textwidth]{figures/arkhangelsk_oblast_monthly_boxplot.png}
\caption{Распределение числа браков по месяцам: Архангельская область}
\label{fig:ark_box}
\end{figure}

В Архангельской области смена сезонных минимумов и максимумов выражена резче,
чем в Москве, если оценивать ее в относительных величинах. Это хорошо согласуется
с меньшей численностью населения, большей зависимостью от климатических условий
и более компактной демографической базой для формирования браков. Таким образом,
сезонность присутствует в обоих регионах, но в северном регионе она проявляется
контрастнее.

\section{HEGY-type тест и порядок интегрируемости}
\subsection{Эконометрическая постановка}
Для месячных рядов важно различать обычные и сезонные единичные корни.
Сезонный единичный корень означает, что шок может сохраняться не только на
уровне долгосрочного тренда, но и в годовой сезонной компоненте с периодом
12 месяцев.

Нулевая гипотеза HEGY-теста: временной ряд содержит сезонные и/или несезонные
единичные корни (ряд нестационарен). Отвержение нулевой гипотезы означает, что
для соответствующей частоты ряд можно считать стационарным. Если нулевая гипотеза
не отвергается для уровня ряда, но отвергается после обычной и сезонной
разностей, это служит основанием для выбора порядка интегрируемости.

В данной работе использовалось приближение HEGY-теста в виде сезонной регрессии
единичного корня с двумя спецификациями: только с константой и с константой,
трендом и месячными фиктивными переменными. Такой подход следует трактовать как
\emph{HEGY-type seasonal unit root regression}, а не как полную официальную
реализацию HEGY.

\subsection{Результаты для Москвы}
В таблице~\ref{tab:hegy_moscow} приведены результаты HEGY-type регрессии для
Москвы.

\begin{table}[H]
\centering
\caption{HEGY-type регрессия единичного корня: Москва}
\label{tab:hegy_moscow}
\small
\begin{tabularx}{\textwidth}{l >{\RaggedRight\arraybackslash}p{4.3cm} l r r l}
\toprule
Ряд & H0 & Спецификация & Статистика & $p$-value & Вывод \\
\midrule
$y_t$ & \makecell[l]{ряд содержит сезонные\\и/или несезонные\\единичные корни} & C & 17.9550 & 0.0000 & нестационарен \\
$y_t$ & \makecell[l]{ряд содержит сезонные\\и/или несезонные\\единичные корни} & C+T+M & 0.0090 & 0.9931 & нестационарен \\
$\Delta y_t$ & \makecell[l]{ряд содержит сезонные\\и/или несезонные\\единичные корни} & C & 12.6100 & 0.0000 & нестационарен \\
$\Delta y_t$ & \makecell[l]{ряд содержит сезонные\\и/или несезонные\\единичные корни} & C+T+M & 0.8470 & 0.3987 & нестационарен \\
$\Delta_{12} y_t$ & \makecell[l]{ряд содержит сезонные\\и/или несезонные\\единичные корни} & C & -4.0410 & 0.0001 & стационарен \\
$\Delta_{12} y_t$ & \makecell[l]{ряд содержит сезонные\\и/или несезонные\\единичные корни} & C+T+M & -4.1950 & 0.0001 & стационарен \\
$\Delta \Delta_{12} y_t$ & \makecell[l]{ряд содержит сезонные\\и/или несезонные\\единичные корни} & C & -2.2740 & 0.0250 & стационарен \\
$\Delta \Delta_{12} y_t$ & \makecell[l]{ряд содержит сезонные\\и/или несезонные\\единичные корни} & C+T+M & -2.1330 & 0.0356 & стационарен \\
\bottomrule
\end{tabularx}
\end{table}

Для Москвы уровень ряда и ряд только с первой разностью сохраняют признаки
нестационарности. Напротив, сезонная разность по лагу 12 и комбинированная
разность дают стационарное представление. Следовательно, для Москвы требуется
одна несезонная и одна сезонная разность, то есть выбор \((d,D)=(1,1)\).

\subsection{Результаты для Архангельской области}
В таблице~\ref{tab:hegy_ark} представлены результаты HEGY-type регрессии для
Архангельской области.

\begin{table}[H]
\centering
\caption{HEGY-type регрессия единичного корня: Архангельская область}
\label{tab:hegy_ark}
\small
\begin{tabularx}{\textwidth}{l >{\RaggedRight\arraybackslash}p{4.3cm} l r r l}
\toprule
Ряд & H0 & Спецификация & Статистика & $p$-value & Вывод \\
\midrule
$y_t$ & \makecell[l]{ряд содержит сезонные\\и/или несезонные\\единичные корни} & C & 21.2400 & 0.0000 & нестационарен \\
$y_t$ & \makecell[l]{ряд содержит сезонные\\и/или несезонные\\единичные корни} & C+T+M & 0.4250 & 0.6720 & нестационарен \\
$\Delta y_t$ & \makecell[l]{ряд содержит сезонные\\и/или несезонные\\единичные корни} & C & 20.4340 & 0.0000 & нестационарен \\
$\Delta y_t$ & \makecell[l]{ряд содержит сезонные\\и/или несезонные\\единичные корни} & C+T+M & 0.1060 & 0.9158 & нестационарен \\
$\Delta_{12} y_t$ & \makecell[l]{ряд содержит сезонные\\и/или несезонные\\единичные корни} & C & -4.1770 & 0.0001 & стационарен \\
$\Delta_{12} y_t$ & \makecell[l]{ряд содержит сезонные\\и/или несезонные\\единичные корни} & C+T+M & -4.9010 & 0.0000 & стационарен \\
$\Delta \Delta_{12} y_t$ & \makecell[l]{ряд содержит сезонные\\и/или несезонные\\единичные корни} & C & -3.5460 & 0.0006 & стационарен \\
$\Delta \Delta_{12} y_t$ & \makecell[l]{ряд содержит сезонные\\и/или несезонные\\единичные корни} & C+T+M & -3.3870 & 0.0010 & стационарен \\
\bottomrule
\end{tabularx}
\par\smallskip
\small Примечание: C --- константа, T --- линейный тренд, M --- месячные
фиктивные переменные.
\end{table}

Для Архангельской области картина аналогична московской. Уровень ряда и первая
разность не обеспечивают стационарности, тогда как сезонная и комбинированная
разности позволяют отвергнуть нестационарность. Тем самым и для второго региона
выбор \((d,D)=(1,1)\) оказывается согласованным как с тестовой статистикой, так
и с визуальной диагностикой.

\subsection{Графическая диагностика}
На рисунке~\ref{fig:moscow_diff} показаны исходный логарифмический ряд Москвы и
его разности.

\begin{figure}[H]
\centering
\includegraphics[width=0.95\textwidth]{figures/moscow_differencing_diagnostics.png}
\caption{Диагностика разностей: Москва}
\label{fig:moscow_diff}
\end{figure}

После применения сначала обычной, а затем сезонной разности колебания московского
ряда становятся заметно более стабильными вокруг постоянного уровня. Это
поддерживает вывод HEGY-type анализа о необходимости одновременно несезонного и
сезонного дифференцирования.

На рисунке~\ref{fig:ark_diff} представлена аналогичная диагностика для
Архангельской области.

\begin{figure}[H]
\centering
\includegraphics[width=0.95\textwidth]{figures/arkhangelsk_oblast_differencing_diagnostics.png}
\caption{Диагностика разностей: Архангельская область}
\label{fig:ark_diff}
\end{figure}

Для Архангельской области после сезонной и комбинированной разностей уровень
вариации также заметно сглаживается. Следовательно, выбор \((d,D)=(1,1)\)
обоснован не только статистически, но и визуально.

\section{SARIMA-модели}
\subsection{Набор сравниваемых спецификаций}
Для Москвы сравнивались четыре мультипликативные SARIMA-спецификации.
Модель SARIMA$(0,1,1)(0,1,1)_{12}$ рассматривалась как базовая сезонная
MA-альтернатива для ряда с выраженными пиками на первом и двенадцатом лагах.
Модель SARIMA$(1,0,0)(1,1,0)_{12}$ включена в отчет как обязательная
AR-альтернатива по формулировке задания и служит жестким benchmark-вариантом без
компонент скользящего среднего. Модель SARIMA$(1,1,1)(1,1,1)_{12}$ выступает
как наиболее гибкий сбалансированный кандидат, а спецификация
SARIMA$(1,1,1)(0,1,1)_{12}$ добавлена на основе наблюдавшихся паттернов ACF/PACF.

Информационные критерии AIC и BIC отражают компромисс между качеством подгонки
и числом параметров: чем ниже значение критерия, тем модель предпочтительнее при
прочих равных. Однако окончательное решение принимается только с учетом
диагностики остатков, а также проверки стационарности и обратимости.

\subsection{Формы ключевых моделей}
Обозначая оператор лага через \(L\), получаем следующие мультипликативные формы:
\begin{align*}
(1 - L)(1 - L^{12}) y_t
  &= (1 + \theta_1 L)(1 + \Theta_1 L^{12}) \varepsilon_t, \\
(1 - \phi_1 L)(1 - \Phi_1 L^{12})(1 - L)(1 - L^{12}) y_t
  &= (1 + \theta_1 L)(1 + \Theta_1 L^{12}) \varepsilon_t.
\end{align*}
Первое равенство соответствует модели SARIMA$(0,1,1)(0,1,1)_{12}$, второе ---
модели SARIMA$(1,1,1)(1,1,1)_{12}$.

\subsection{Сравнение моделей для Москвы}
В таблице~\ref{tab:moscow_models} приведено сравнение моделей для Москвы.

\begin{table}[H]
\centering
\caption{Сравнение SARIMA-моделей для Москвы}
\label{tab:moscow_models}
\small
\begin{tabularx}{\textwidth}{>{\RaggedRight\arraybackslash}p{3.3cm}
                            r
                            r
                            >{\RaggedRight\arraybackslash}p{1.8cm}
                            >{\RaggedRight\arraybackslash}p{1.7cm}
                            Y}
\toprule
Модель & AIC & BIC & Диагностика & Корни & Комментарий \\
\midrule
\makecell[l]{SARIMA$(0,1,1)$\\$(0,1,1)_{12}$} & -102.78 & -94.79 & пройдена & нет &
Сильная MA-структура по обычному и сезонному лагу, но условие обратимости
не-сезонной MA-части не выполняется. \\
\makecell[l]{SARIMA$(1,0,0)$\\$(1,1,0)_{12}$} & -37.81 & -29.79 & не пройдена & да &
Жесткая AR-альтернатива, включенная в соответствии с заданием; по качеству
подгонки и остаточной диагностике она заметно уступает более гибким моделям. \\
\makecell[l]{SARIMA$(1,1,1)$\\$(1,1,1)_{12}$} & -102.14 & -88.82 & пройдена & да &
Наиболее сбалансированная спецификация: хорошие AIC/BIC, приемлемые остатки и
корректные корневые диагностики. \\
\makecell[l]{SARIMA$(1,1,1)$\\$(0,1,1)_{12}$} & -108.71 & -98.06 & не пройдена & да &
Минимальный AIC указывает на хорошую подгонку, но остаточная зависимость остается
слишком заметной для итогового выбора. \\
\bottomrule
\end{tabularx}
\end{table}

Хотя по AIC минимальна модель SARIMA$(1,1,1)(0,1,1)_{12}$, она не проходит
диагностику остатков. AR-альтернатива SARIMA$(1,0,0)(1,1,0)_{12}$ также
оказывается слабой: ее AIC/BIC существенно хуже, а остаточная автокорреляция
сохраняется. Поэтому итоговый выбор делается в пользу
SARIMA$(1,1,1)(1,1,1)_{12}$, которая сочетает низкие значения информационных
критериев с приемлемой диагностикой и выполнением условий стационарности и
обратимости.

\subsection{Итоговая модель для Архангельской области}
В таблице~\ref{tab:ark_model} приведена итоговая модель для Архангельской
области.

\begin{table}[H]
\centering
\caption{Итоговая SARIMA-модель для Архангельской области}
\label{tab:ark_model}
\small
\begin{tabularx}{\textwidth}{>{\RaggedRight\arraybackslash}p{3.3cm}
                            r
                            r
                            >{\RaggedRight\arraybackslash}p{1.8cm}
                            >{\RaggedRight\arraybackslash}p{1.7cm}
                            Y}
\toprule
Модель & AIC & BIC & Диагностика & Корни & Комментарий \\
\midrule
\makecell[l]{SARIMA$(0,1,1)$\\$(0,1,1)_{12}$} & -100.25 & -92.26 & пройдена & да &
Компактная сезонная MA-спецификация с приемлемой остаточной диагностикой и
наиболее убедительным балансом простоты и качества. \\
\bottomrule
\end{tabularx}
\end{table}

Для Архангельской области предпочтение отдано более простой модели
SARIMA$(0,1,1)(0,1,1)_{12}$. Для менее масштабного ряда такая компактная
спецификация оказывается предпочтительнее, поскольку снижает риск избыточной
параметризации и при этом сохраняет содержательно приемлемые диагностические
свойства.

\subsection{Диагностика остатков}
На рисунке~\ref{fig:moscow_resid} показаны остатки и автокорреляционная функция
остатков выбранной модели для Москвы.

\begin{figure}[H]
\centering
\includegraphics[width=0.88\textwidth]{figures/moscow_11111112_residuals.png}
\caption{Остатки и ACF остатков выбранной модели для Москвы}
\label{fig:moscow_resid}
\end{figure}

Для Москвы остатки в целом близки к белому шуму, хотя на коротких лагах
сохраняются слабые следы структуры. Эти отклонения не меняют общего вывода о
приемлемости модели, поскольку выбранная спецификация проходит итоговую
диагностику и одновременно сохраняет устойчивость по корневым критериям.

На рисунке~\ref{fig:ark_resid} представлена аналогичная диагностика для
Архангельской области.

\begin{figure}[H]
\centering
\includegraphics[width=0.88\textwidth]{figures/arkhangelsk_oblast_01101112_residuals.png}
\caption{Остатки и ACF остатков выбранной модели для Архангельской области}
\label{fig:ark_resid}
\end{figure}

Для Архангельской области остаточное поведение выглядит еще более компактным.
Это согласуется с выбором более простой итоговой модели, которая достаточно
хорошо описывает динамику ряда без лишнего усложнения параметризации.

\section{Условие стационарности}
Для выбранной московской модели SARIMA$(1,1,1)(1,1,1)_{12}$ авторегрессионная
характеристика имеет вид
\[
(1 - \phi_1 z)(1 - \Phi_1 z^{12}) = 0.
\]
Условие стационарности требует, чтобы все корни характеристического уравнения
авторегрессионной части лежали вне единичной окружности. Диагностика корней
оцененной модели показывает, что это условие выполняется. Следовательно,
выбранная спецификация удовлетворяет требованию стационарности и может
использоваться для интерпретации и прогноза.

\section{Прогноз}
\subsection{Прогноз для Москвы}
На рисунке~\ref{fig:moscow_forecast} представлен 12-месячный прогноз для Москвы
вместе с 95\%-ным доверительным интервалом.

\begin{figure}[H]
\centering
\includegraphics[width=0.9\textwidth]{figures/moscow_forecast.png}
\caption{12-месячный прогноз для Москвы}
\label{fig:moscow_forecast}
\end{figure}

Для Москвы прогноз сохраняет ту же годовую сезонную волну, что и исторический
ряд: после зимних минимумов ожидается рост к летне-осеннему максимуму, а затем
возврат к более низким уровням. Такая траектория выглядит реалистично для
мегаполиса, где сезонность устойчива, но не сопровождается резкими структурными
сдвигами. Доверительный интервал постепенно расширяется по мере удаления от
последнего наблюдения, что естественно отражает рост прогнозной неопределенности.

\subsection{Прогноз для Архангельской области}
На рисунке~\ref{fig:ark_forecast} показан 12-месячный прогноз для Архангельской
области.

\begin{figure}[H]
\centering
\includegraphics[width=0.9\textwidth]{figures/arkhangelsk_oblast_forecast.png}
\caption{12-месячный прогноз для Архангельской области}
\label{fig:ark_forecast}
\end{figure}

Прогноз для Архангельской области также сохраняет сезонную структуру, но с более
заметной относительной амплитудой колебаний. Это соответствует меньшему размеру
ряда и большей зависимости от климатических и демографических факторов.
Интервалы неопределенности расширяются к концу горизонта, но не указывают на
взрывное поведение или потерю правдоподобия прогноза.

\subsection{Прогнозные значения}
В таблице~\ref{tab:forecast_moscow} приведен 12-месячный прогноз для Москвы.

\begin{table}[H]
\centering
\caption{12-месячный прогноз для Москвы}
\label{tab:forecast_moscow}
\small
\begin{tabular}{lrrr}
\toprule
Период & Точечный прогноз & Нижняя граница 95\% ДИ & Верхняя граница 95\% ДИ \\
\midrule
2017-02 & 4662.60 & 3584.19 & 6065.49 \\
2017-03 & 4550.88 & 3496.02 & 5924.01 \\
2017-04 & 6243.11 & 4761.70 & 8185.40 \\
2017-05 & 4129.02 & 3140.82 & 5428.13 \\
2017-06 & 9057.26 & 6864.21 & 11950.97 \\
2017-07 & 10107.87 & 7634.96 & 13381.73 \\
2017-08 & 10811.13 & 8138.63 & 14361.20 \\
2017-09 & 9691.82 & 7271.86 & 12917.10 \\
2017-10 & 7292.26 & 5453.54 & 9750.91 \\
2017-11 & 5728.94 & 4270.45 & 7685.54 \\
2017-12 & 5796.48 & 4307.33 & 7800.45 \\
2018-01 & 3965.67 & 2937.56 & 5353.61 \\
\bottomrule
\end{tabular}
\end{table}

Для Москвы прогноз указывает на рост от зимних уровней к летнему пику свыше
10 тыс. браков в месяц, после чего ожидается возврат к более низким осенне-зимним
значениям. Такой профиль полностью согласуется с ранее выявленной сезонностью и
не содержит признаков нереалистичного поведения.

В таблице~\ref{tab:forecast_ark} приведен 12-месячный прогноз для Архангельской
области.

\begin{table}[H]
\centering
\caption{12-месячный прогноз для Архангельской области}
\label{tab:forecast_ark}
\small
\begin{tabular}{lrrr}
\toprule
Период & Точечный прогноз & Нижняя граница 95\% ДИ & Верхняя граница 95\% ДИ \\
\midrule
2017-02 & 442.10 & 334.75 & 583.88 \\
2017-03 & 487.90 & 367.58 & 647.60 \\
2017-04 & 502.95 & 377.05 & 670.87 \\
2017-05 & 241.49 & 180.17 & 323.69 \\
2017-06 & 742.68 & 551.45 & 1000.22 \\
2017-07 & 900.19 & 665.28 & 1218.05 \\
2017-08 & 942.16 & 693.09 & 1280.75 \\
2017-09 & 716.37 & 524.59 & 978.26 \\
2017-10 & 539.95 & 393.63 & 740.66 \\
2017-11 & 510.39 & 370.44 & 703.21 \\
2017-12 & 602.76 & 435.58 & 834.12 \\
2018-01 & 335.82 & 241.63 & 466.72 \\
\bottomrule
\end{tabular}
\end{table}

Для Архангельской области прогнозная траектория аналогична по форме, но заметно
ниже по масштабу: от нескольких сотен регистраций в низкий сезон до примерно
900--950 в месяцы пика. Это согласуется с меньшей численностью населения и
более сильной ролью сезонных ограничений.

В целом прогноз отвечает на вопрос о дальнейшей динамике показателя в двух
регионах достаточно однозначно. Для Москвы ожидается сохранение более высокого
уровня брачности и той же сезонной волны с летне-осенним максимумом, тогда как
для Архангельской области прогноз также воспроизводит сезонный профиль, но на
существенно более низком масштабе. В обоих случаях доверительные интервалы
постепенно расширяются по мере удаления от конца выборки, что естественно для
краткосрочного прогноза; при этом траектории не выглядят взрывными или
неправдоподобными, поэтому прогноз можно считать реалистичным.

\section{Заключение}
Проведенный анализ показал, что в обоих регионах присутствует выраженная годовая
сезонность, однако уровень и относительная вариативность существенно различаются.
Москва характеризуется более высоким и устойчивым числом браков, тогда как
Архангельская область сильнее реагирует на сезонные и демографические факторы.

HEGY-type анализ привел к выбору порядка интегрируемости \((d,D)=(1,1)\), то
есть одной несезонной и одной сезонной разности. Для Москвы лучшей оказалась
модель SARIMA$(1,1,1)(1,1,1)_{12}$, а для Архангельской области ---
SARIMA$(0,1,1)(0,1,1)_{12}$. Построенные прогнозы сохраняют сезонный профиль, а
доверительные интервалы расширяются плавно, что делает итоговые прогнозы
реалистичными и пригодными для краткосрочного анализа.

\end{document}
